There was no evacuating a city from the inside of itself. That was the thing Davos had spent four hours failing to explain to men who kept coming to him for orders, and had spent the fifth hour giving up on explaining and simply doing what could be done instead.

An evacuation had a direction. That was the whole of what made it an evacuation — you had a bad place and a better one and you moved people from the first toward the second, and the entire art of it was working out the roads and the rate and how many you could push down a given street in an hour without killing forty of them in a crush. He had planned three of those in his life and he had planned this one for four days with Tyrion’s chalk marks and a map of the city and every route out marked in three colours, and all of it had been built on the assumption that the danger was north of the walls.

The danger was under the fishmonger’s. The danger was in the crypt stair of every poor sept in Flea Bottom and in four hundred years of cellars and in a drain outside a cobbler’s shop in Cobbler’s Square, and there was no direction to move anybody, because every street in King’s Landing had the same thing under it.

So he had thrown the map away at the second hour and started again with the only principle left to him, which was \textit{up}.

\sceneBreak

“Rope-walk lofts,” he said. “Granary lofts. Any upper floor with one stair and no cellar. Bell towers. The Dragonpit’s ruined but the seating tiers are stone and there’s nothing under them but rubble — put people in the tiers and hold the two approaches with twenty men instead of holding forty streets with none.” He had his sergeants in a half circle in the courtyard of a livery stable and he was drawing on the flags with a piece of charcoal because there was no time for paper. “Don’t send them anywhere. Sending them anywhere is what’s killing them — every crowd on the move is a crowd going past twelve cellars. Get them off the ground where they already are and \textit{leave them there} and mark the door, and we’ll come back for them when there’s a coming back.”

“There’s people won’t go up, my lord. They want out. They keep saying the north gate — “

“I know what they keep saying.” Davos straightened. “Tell them the truth. Tell them there’s ninety thousand of the things outside the north gate and about nine hundred of us between them and it, and that the only reason any of us is still breathing is that they haven’t got through the trench yet. Tell them a loft in Pisswater is the safest place in this city and that I said so and that I’m not going anywhere myself.” He looked round at them. “And if they won’t have it from you, get a septon. There’s four of them working the Street of Sisters and people will take from a septon what they’ll not take from a man with a sword.”

They went. Davos stood a moment in the empty courtyard with charcoal on his fingers and let himself feel, for about the length of a breath, the full size of what he was doing arithmetic on.

He had about six hundred men that anybody would call soldiers and could not spare more than two hundred of them from the works. There were somewhere north of two hundred thousand people in King’s Landing. There were, by the best guess anyone had — Bronn’s guess, which Davos had checked twice because he had not wanted it to be Bronn’s — perhaps thirty thousand bodies in the ground under the city, in pits and cellars and crypts, of which nobody had the faintest idea how many were coming up.

\textit{You cannot save two hundred thousand people with two hundred men.}

That was true. He had known it since the second hour. What he had also worked out, somewhere in the fourth, was that this was not actually the problem in front of him — because he was not going to save two hundred thousand people, and knowing it changed nothing about what to do next. The problem in front of him was a street. Then another street. You did the street you were standing in and you did not permit yourself the luxury of the whole sum, because the whole sum would stop your hands, and hands were the only thing anybody had left.

He had learned that on a burning river twenty years ago. It had taken him until tonight to find out he still knew it.

\sceneBreak

Varys found him at the Street of Steel at the sixth hour, and Davos was so glad to see him that it embarrassed him.

“Four hundred in the Guild of Alchemists’ hall,” Varys said, without preamble, which was itself a measure of the night. “It’s stone, the cellar’s been bricked up since Aerys and I have confirmed that personally rather than taking a guildsman’s word for it. Another three hundred or so in the tiers of the Dragonpit, and rising — your notion about the tiers was the best thing anyone’s said today, Lord Davos, and I’ve stolen it four times without crediting you.”

“Steal it forty times.” Davos wiped his face. “Where are you getting your numbers?”

“The same place I have always got my numbers.” Something almost like a smile. “There are a great many people in this city who have carried messages for me across twenty years, and it transpires that a woman who will carry a message will also carry a count, and that people who are frightened will tell a familiar face things they would not tell a lord with a sword. My little birds have all grown up, Lord Davos, and rather a lot of them have children of their own, and this evening every one of them is telling me which cellars are open and which streets are clear.” He said it lightly and then, for a moment, did not. “I have spent my life being told that what I do is not honest work. I find I have never in my life been so glad to have done it.”

They stood together a moment in the ruin of the Street of Steel while somebody’s forge, unattended for six hours, went on glowing behind them.

“You’ll want to be careful,” Davos said. “You’re going down into places tonight that a man like you has no business in.”

“A man like me.” Varys considered that. “I was born in a place considerably worse than any cellar in this city, and I was sold out of it, and I have spent every year since making certain that I would never again be somewhere I could not get out of.” He adjusted his sleeves, an old fussy gesture that Davos had seen a hundred times in council chambers and had never once seen mean anything. “I find that tonight I am going down into them anyway, and that this is the first thing I have done in forty years that I would not need to explain to anybody. It’s rather restful. Do go on being careful yourself. Someone has to be left who can do the arithmetic.”

\sceneBreak

They brought Davos word at the seventh hour of survivors found in a rope-walk loft off Pisswater Bend. Twelve of them — eleven and a girl of about nine — pulled up a ladder into a room with one stair and no cellar under it, alive, all of them, in a street where his sergeants counted forty dead on the paving below.

“Who got them up there?” Davos asked.

“Nobody knows, my lord. They said it was a woman. Big woman, older, carried a bucket.” The sergeant shrugged, not unkindly; he had eleven other things to report. “She went back down for another one and didn’t come up again. They don’t have a name for her. One of them said she hit her.”

Davos wrote \textit{12 saved, rope-walk loft, Pisswater} in the tally he was keeping on the back of Tyrion’s ruined chalk map, because he was keeping a tally and had decided at the third hour that he would keep it honestly.

He did not have anywhere to write the rest of it. He thought about that for a moment and then went on to the next street.

\sceneBreak

He found Qyburn at the ninth hour, in the lower levels of the Red Keep, and did not at first believe what he was looking at.

They had turned the old kitchens into a hospital because the kitchens were the largest room in the castle with no cellar beneath them, and there were perhaps two hundred people on the flags in there, and in the middle of it, in shirtsleeves with his forearms wet to the elbow, the former Grand Maester of the Seven Kingdoms was setting a woman’s leg.

Davos stood in the doorway for a long time.

He was \textit{good}. That was the thing. Davos had watched maesters work and had watched battlefield surgeons work and had once, on a boat, watched a man saw off another man’s foot with a carpenter’s tool, and he knew the difference between someone doing his best and someone who simply \textit{knew}, and Qyburn knew. He worked fast and he did not waste motion and he did not lie to people about what was about to happen to them. He had four girls from the city working under him and he had plainly taught them something in the last few hours, because they anticipated him. Twice while Davos watched, he stopped what he was doing to walk down the line and reorder it — sending one man to the back who had been brought in first, moving a boy up who had been lying quiet in a corner because he was too polite to complain, and Davos, who had spent the whole night doing exactly that kind of arithmetic, recognized it instantly and felt something in his chest that he did not want to be feeling.

“How long has he been here?” he asked one of the girls.

“Since before dark, my lord. He come down on his own. Nobody sent for him.” She was fifteen at the outside and had blood to her wrists. “He don’t sleep and he don’t eat and he don’t talk much except to tell you what to do. He fixed my uncle’s arm.” She glanced across at him. “They say he done things. Up in the tower. Before.”

“Aye,” Davos said. “They do.”

He should have gone to find Tyrion. He knew that, standing there. Tyrion had put out an order four days ago to find Qyburn, and Qyburn had a great deal to answer for, and the correct thing — the orderly thing, the thing a man on a council does — was to send two men to take him into custody and to keep him alive for whatever came after, because the man knew the wildfire signal and knew this castle and was, in the coldest possible sense, an asset.

Davos did not go and find Tyrion. He would spend the rest of his life being unable to say with certainty why not.

He told himself, at the time, that it was because there were two hundred people on that floor and one man who could work on them and no one to spare who could do it half so well. That was true. It was true and it was not the whole of it, and he had been alive too long to pretend otherwise even to himself.

He stayed for three hours and helped carry.

\sceneBreak

The woman came in a little after midnight, and Davos knew her, which was the part that ruined it.

Her name was Wyla and she sold — had sold — cured fish out of a barrow on the Hook, and she had been working beside Davos since the fourth hour. She had gone down into two cellars that he personally would not have gone into. She had carried a child up a ladder with the child’s mother screaming at her from below for being too slow, and had said nothing about it afterward. She had brought him three separate corrections to his own count of who was where, and every one of them had been right. Davos had been intending, when there was time, to find out whether she wanted work.

She came into the kitchens with a man’s arm over her shoulder and she got him down onto the flags and she straightened up, and she saw Qyburn across the room, and she stopped moving entirely.

“That’s him,” she said.

Nobody heard her the first time. She said it again, louder, and this time about forty people heard her.

“That’s him. That’s the one from the tower.” Her voice had gone strange and flat. “That’s the one. He took my Marlon. Ten years, near enough. He took my Marlon out of the black cells and I went up there and I asked and asked, and they told me he’d been let go, they told me he’d been \textit{let go}, and I’ve never — “ She stopped. She started again and it came out worse. “I asked him to his face. I stood in front of that man and asked him to his face and he said he had no memory of the name.”

The room had gone quiet in the way a room goes quiet.

Qyburn straightened up from what he was doing.

Davos watched him look at her. He watched, very carefully, because he wanted to know — and what he saw was a man searching a considerable and well-ordered memory, honestly, with no fear in his face at all, and finding the name, and \textit{placing it}.

“Marlon,” Qyburn said. “Yes. A thin lad. Fifteen or sixteen.” He said it in the same tone he had used all night to tell people what was about to happen to them: informatively, without malice, without apology. “He was of considerable use to me. I learned a great deal from him about the sequence in which a body fails, which I have applied several times tonight in this room, including twice to men who are presently alive because of it.”

He did not say it to be cruel. Davos understood that, then and afterward and for the rest of his life. The man was not taunting her. He was offering her, in complete sincerity, what he believed to be the relevant information — an \textit{accounting}, a set of entries on either side of a line — and he had genuinely no idea, none at all, that he had just said the worst sentence available in the language.

Somebody threw the first thing. Davos never saw what it was or who.

\sceneBreak

There were perhaps sixty people in that room who could stand.

It was not fast. He would have preferred, afterward, to be able to say it was fast. There was a great deal of noise and at some point Qyburn went down and after that Davos could not see him at all, only the backs of people, and the noise changed twice — once when the shouting stopped being words, and once, later, when it stopped being much of anything.

He heard the man’s voice twice through it. Both times it was the same thing, or near enough, and both times it had no fear in it and no pleading either — only a kind of appalled bureaucratic bewilderment, the voice of someone being told his papers are not in order:

“I have saved \textit{forty men in this room tonight} — “

And then, later, nearly lost in it:

“\textit{This is not proportionate.}”

Davos stood eight feet away with his back against the wall of the old kitchens and his hands at his sides, and did not move.

\sceneBreak

He could have stopped it. That was the thing he would take out and look at, afterward, in the small hours, for whatever years he had left.

Not easily. But he was Davos Seaworth, Hand to a dead king and a councillor of this one, and he was known by sight to half that room, and if he had put himself between Wyla and that man and used the voice he used on sergeants, then some of them would have stopped, and the ones who stopped would have stopped the rest. It would have cost him something. It might have cost him a broken arm and it would certainly have cost him Wyla, who had gone down into two cellars for him tonight and who would never afterward have looked at him with anything but contempt.

He had stood in front of a red woman once, in a yard at Winterfell, and told her plainly that he did not forgive her for burning a child and would not pretend to, and that he would not stop her from what she meant to do next either, because what she meant to do next was useful and he was not fool enough to pretend that usefulness and forgiveness were the same coin.

He had hated himself for weeks over that. He had told himself, in the worst of those weeks, that he had chosen the war over the girl, and that a better man would have chosen otherwise.

And here, tonight, in a room full of people who had been sold and used and taken out of black cells and told their sons had been \textit{let go} — here he had gone the other way. He had let usefulness go, and let the account be settled by the only people who had ever had any standing to settle it.

He did not know whether that was the same man learning something, or a different man giving himself permission.

He genuinely did not know. That was the whole of it. He would turn it over for years and it would never come out the same way twice, and he came to understand, eventually, that some sums are not meant to close, and that a person who insists on closing them is generally lying about at least one of the entries.

\sceneBreak

The four girls he had taught went back to work about a quarter hour later, because there were two hundred people on that floor and nobody else to do it.

They were not as good as he had been. They were considerably better than nothing. One of them — the fifteen-year-old with blood to her wrists, whose uncle’s arm he had fixed — reordered the line the way she had watched him do it, sending a man to the back and moving a quiet boy up, and got it right.

Davos noticed that, and noticed that nobody else in the room had.

He did not write it down. He had a tally in his pocket on the back of a chalk map and it had \textit{12 saved, rope-walk loft, Pisswater} on it in his own laborious hand, and he thought for a moment about what he would even put — and found there was no column for it, and that if he made one, the only honest entry would have to have the rest of it beside it, and that he was not prepared to write that down in a hand that other people would one day read.

So there was no record.

He went back out into the city and did the next street.